\documentclass[../main.tex]{subfiles}
\graphicspath{{\subfix{../img/}}}
\begin{document}
	\subsection{Carica elettrone}
	La carica di un elettrone ($e^-$) è: $-1.6 \cdot 10^{-19} C$ \\
	L'unita di misura utilizzata è il Coulomb.
	
	\subsection{Prima legge di Ohm}
	\begin{equation}
		\Delta V = R \cdot I \;\; [V]
	\end{equation}
	$\Delta V$ è la \textbf{differenza di potenziale} che si misura in Volt $[V]$ \\
	$R$ è la \textbf{resistenza} che si misura in Ohm $[\Omega]$ \\
	$I$ è la \textbf{corrente elettrica} che si misura in Ampere $[A]$\\\\
	La differenza di potenziale \textbf{può anche essere negativa}, ma in quel caso indica che cambia il verso.
	\subsection{Seconda legge di Ohm}
	\begin{equation}
		R = \rho \cdot \frac{l}{S} \;\; [\Omega \cdot m]
	\end{equation}
	$R$ indica la resistenza \\
	$\rho$ indica la \textbf{resistività}, ovvero quanto un materiale resiste all'elettricità, $\rho \in \left[ 10^{-7}, 10^{16}\right]$ \\
	$l$ indica la \textbf{lunghezza} del campione in cui scorre la corrente \\
	$S$indica la \textbf{sezione/area} del campione in cui scorre la corrente \\ \\
	L'inverso della resistività è la \textbf{conducibilità}, ovvero quanto un materiale conduce elettricità, viene indicata con:
	\[ \sigma = \frac{1}{\rho} \]
	\subsubsection{Variazione di Resistività (o Concudibilità) con la temperatura}
	\begin{equation}
		\rho = \rho_0 \cdot \left[ 1 + \alpha \cdot (T - T_0)\right]
	\end{equation}
	\begin{itemize}
		\item $\rho_0$ è la resistività calcolata ad una certa temperatura $T_0$
		\item $\alpha$ è il coefficiente termico di resistività che cambia a seconda del materiale, si musira in $\left[ \frac{1}{ ^\circ C}\right]$
	\end{itemize}
	Un aumento di temperatura spesso comporta un aumento di resistività.
	\subsection{Bipoli}
	\subsubsection{Bipolo Elettrico (generico)}
	\begin{center}
		\begin{figure}[h]
			\centering
			\includegraphics[scale=0.4]{1.jpeg}
			\caption{Bipolo Elettrico}
			\label{fig:im-1}
		\end{figure}
	\end{center}
	\subsubsection{Filo Elettrico}
	\subsubsection{Bipolo "Resistenza"}
	\begin{center}
		\begin{figure}[h]
			\centering
			\includegraphics[scale=0.4]{2.png}
			\caption{Bipolo Resistenza}
			\label{fig:im-2}
		\end{figure}
	\end{center}
	\subsubsection{Bipolo "Generatore Ideale di Tensione"}
	\begin{figure}
		\centering
		\begin{subfigure}{.5\textwidth}
			\centering
			\includegraphics[width=.5\linewidth]{3.jpeg}
			\caption{Bipolo Generatore Ideale di Tensione}
			\label{fig:im-3a}
		\end{subfigure}%
		\begin{subfigure}{.5\textwidth}
			\centering
			\includegraphics[width=.5\linewidth]{3b.jpeg}
			\caption{Bipolo Generatore Ideale di Tensione}
			\label{fig:im-3b}
		\end{subfigure}
		\caption{Bipolo Generatore Ideale di Tensione}
		\label{fig:im-3}
	\end{figure}
	\subsubsection{Bipolo "Generatore Ideale di Corrente"}
	\begin{center}
		\begin{figure}[h]
			\centering
			\includegraphics[scale=0.1]{4.jpeg}
			\caption{Bipolo Generatore Ideale di Corrente}
			\label{fig:im-4}
		\end{figure}
	\end{center}
	\subsection{Convenzione del Generatore}
	Convenzione in cui la tensione $\Delta V$ e la corrente $I$ elettrica hanno lo stesso verso, in questo caso la prima legge di Ohm diventa:
	\begin{equation}
		\Delta V = - (R \cdot I)
	\end{equation}
	\subsection{Convenzione del'Utilizzatore}
	Convenzione in cui la tensione $\Delta V$ e la corrente $I$ elettrica hanno verso discorde, in questo caso la prima legge di Ohm diventa:
	\begin{equation}
		\Delta V = R \cdot I
	\end{equation}
	\subsection{Principio di Kirchhoff al Nodo (1° Principio di Kirchhoff)}
	Il primo principio di Kirchhoff afferma che, dato un nodo in cui convergono 3 o più fili, ciascuno con una propria corrente $I_1,I_2, \dots,I_n$, la sommatoria di tutte le correnti entranti è nulla, ovvero la sommatoria delle correnti entranti è uguale alla sommatoria delle correnti uscenti dal nodo.
	\subsection{Principio di Kirchhoff alla Maglia (2° Principio di Kirchhoff)}
	Il secondo principio di Kirchhoff afferma che,data una maglia, la sommatoria ordinata (cioè secondo un unico verso, orario o antiorario, del percorso) di tutte le tensioni della maglia è nulla.
	\subsection{Conduttanza Elettrica}
	\begin{equation}
		G = \frac{1}{R} \;\; [\Omega ^{-1}] = [S]
	\end{equation}
	L'unità di misura della conduttanza è il Siemens
	\subsection{Potenza Elettrica di un Bipolo}
	L'unità di misura della potenza elettrica è il Watt
	\subsubsection{Convenzione Utilizzatore}
	\begin{equation}
		P_{Assorbita} = \Delta V \cdot I \;\; [W]
	\end{equation}
	\begin{equation}
		P_{Generata} = -\Delta V \cdot I \;\; [W]
	\end{equation}
	\subsubsection{Convenzione Generatore}
	\begin{equation}
		P_{Assorbita} = -\Delta V \cdot I \;\; [W]
	\end{equation}
	\begin{equation}
		P_{Generata} = \Delta V \cdot I \;\; [W]
	\end{equation}
	\subsubsection{Potenza Assorbita di una Resistenza}
	È sempre positiva:
	\begin{equation}
		P_{Assorb} = \Delta V \cdot I = R \cdot I \cdot I = R \cdot I^2 
	\end{equation}
	\begin{equation}
		P_{Assorb} = \Delta V \cdot I = \Delta V \cdot \frac{\Delta V}{R} = \frac{(\Delta V)^2}{R}
	\end{equation}
	\subsection{Fenomeni Termici}
	\subsubsection{Legge di Ohm Termica}
	\begin{equation}
		T_{Fin} = T_{Amb} + P_{Assorbita} \cdot R_T \;\; [^\circ C]
	\end{equation}
	Dove $R_T$ è la Resistenza Termica
	\subsubsection{Resistenza Termica}
	\begin{equation}
		R_T = \frac{T_{Fin} - T_{Amb}}{P_{Assorb}} \;\; \left[ \frac{^\circ C}{W}\right]
	\end{equation}
	\subsection{Collegamenti fra Bipoli}
	\subsubsection{Collegamento in Serie}
	In una serie di Bipoli, tutte le correnti sono uguali fra loro:
	\begin{equation}
		I_1 = I_2 = \dots = I_n
	\end{equation}
	Mentre invece utilizzando il secondo principio di Kirchhoff otteniamo che la tensione della serie è uguale alla somma delle tensioni di ogni bipolo:
	\begin{equation}
		\Delta V = \sum_{i = 1}^{n} \Delta V_i
	\end{equation}
	\textbf{Una serie di Bipoli è essa stessa un bipolo.}
	\subsubsection{Collegamento in Parallelo}
	In un parallelo di Bipoli, tutte le tensioni sono uguali fra loro:
	\begin{equation}
		\Delta V_1 = \Delta V_2 = \dots = \Delta V_n
	\end{equation}
	Mentre invece utilizzando il primo principio di Kirchhoff otteniamo che la corrente della serie è uguale alla somma delle correnti di ogni bipolo:
	\begin{equation}
		I = \sum_{i = 1}^{n} I_i
	\end{equation}
	\subsubsection{Partitore di Tensione (Serie di Resistenze)}
	\begin{equation}
		R_{Eq} = R_1 + R_2 + \dots + R_n
	\end{equation}
	\begin{equation}
		V_{i} = V_{AB} \cdot \frac{R_i}{\sum_{j=1}^{n} R_j} 
	\end{equation}
	Dove $A$ e $B$ sono i capi della serie di resistori e dove $R_i$ e $V_i$ sono rispettivamente la resistenza e la tensione che del resistore/bipolo preso in considerazione.\\
	Ad esempio, se avessimo due resistenze, $R_n = R_2$
	\subsubsection{Partitore di Corrente (Parallelo di Resistenze)}
	\begin{equation}
		R_{Eq} = \frac{1}{R_1} + \frac{1}{R_2} + \dots + \frac{1}{R_n}
	\end{equation}
	\begin{equation}
		I_i = I_{AB} \cdot \frac{G_i}{\sum_{j = 1}^{n} G_j} = I_{AB} \cdot \frac{\frac{1}{R_i}}{\sum_{j=1}^{n} \frac{1}{R_j}}
	\end{equation}
	\subsubsection{Serie di Generatori Ideali di Tensione}
	\begin{equation}
		E_{Eq} = \sum_{i = 1}^{n} E_i
	\end{equation}
	\subsubsection{Parallelo di Generatori Ideali di Tensione}
	{\color{red} Non esiste è un ASSURDO!!}
	\subsubsection{Serie di Generatori Ideali di Corrente}
	{\color{red} Non esiste è un ASSURDO!!}
	\subsubsection{Parallelo di Generatori Ideali di Corrente}
	\begin{equation}
		I_{Eq} = \sum_{i = 1}^{n} I_i
	\end{equation}
	\subsection{Generatori Reali}
	\subsubsection{Generatore Reale di Tensione}
	È composto da una resistenza $R$, la cui tensione è $V_R = R \cdot I$, ed un generatore ideale di tensione $E_0$, collegati in serie.\\
	La sua legge costitutiva è:
	\begin{equation}
		V_{AB} = E_0 - V_R = E_0 - R \cdot I
	\end{equation}	
	da cui si ricava che $E_0$ ha verso opposto rispetto a $V_{AB}$ e $V_R$.\\
	In un grafico corrente-tensione, la legge è descritta da una retta la cui pendenza è $R$.
	\begin{itemize}
		\item $V_{AB} = E_0$ tensione a vuoto quando $I=0$: Circuito Aperto
		\item $I = \frac{E_0}{R}$ Corrente di Corto Circuito ($I_{CC}$)
	\end{itemize}
	\subsubsection{Generatore Reale di Corrente}
	È composto da una resistenza $R$, la cui corrente è $I_R = \frac{V_{AB}}{R}$, ed un generatore ideale di corrente $I_0$, collegati in parallelo.\\
	La sua legge costitutiva è:
	\begin{equation}
		I = I_0 + I_R = I_0 - \frac{V_{AB}}{R}
	\end{equation}
	In un grafico corrente-tensione, la legge è descritta da una retta la cui pendenza è $R$.\\
	\begin{itemize}
		\item $V_{AB} = R \cdot I_0$ Tensione a vuoto: Circuito Aperto
		\item $I = I_0$ Corrente di Cortocircuito
	\end{itemize}
	\subsubsection{Equivalenze fra Generatori Reali}
	È sempre possibile trasformare un generatore reale di tensione in uno reale di corrente, e viceversa, a patto che vengano rispettate le condizioni di equivalenza:
	\begin{equation}
		\begin{cases}
			E_0 = R \cdot I_0 \\
			I_0 = \frac{E_0}{R}
		\end{cases}
	\end{equation}
	\subsubsection{Considerazioni Energetiche sui Generatori Reali}
	L’equivalenza tra i 2 generatori è solamente elettrica, non  energetica, infatti la potenza generata dai due può variare a seconda del circuito.\\
	Quindi nel caso si debba calcolare la potenza di un circuito è sconsigliato l'utilizzo delle equivalenze.
	\subsubsection{Generatori Reali di Tensione in Serie}
	Formato da $n$ generatori reali di tensione in serie.
	\begin{equation}
		\begin{cases}
			E_{Eq} = E_1 + E_2 + \dots + E_n \\
			R_{Eq} = R_1 + R_2 + \dots + R_n 
		\end{cases}
	\end{equation}
	Se uno dei generatori ha verso invertito rispetto al senso scelto, allora la sua tensione è negativa.
	\subsubsection{Generatori Reali di Corrente in (Serie) Parallelo}
	Formato da $n$ generatori reali di corrente in parallelo.
	\begin{equation}
		\begin{cases}
			I_{Eq} = I_1 + I_2 + \dots + I_n \\
			G_{Eq} = G_1 + G_2 + \dots + G_n = \frac{1}{R_1} + \dots + \frac{1}{R_n}
		\end{cases}
	\end{equation}
	\subsection{Strumenti di Misura}
	\subsubsection{Voltmetro}
	Strumento utilizzato per misurare la tensione di un circuito.\\
	Viene sempre messo in parallelo al circuito in cui si vuole eseguire la misurazione.
	\subsubsection{Amperometro}
	Strumento utilizzato per misurare la corrente di un circuito.\\
	Viene sempre messo in serie al circuito in cui si vuole eseguire la misurazione.
	\subsubsection{Ohmetro}
	Strumento utilizzato per misurare una resistenza.\\
	Spesso non preciso quando si misuarno piccole resistenze.\\
	Viene sempre messo ai capi del resistore su cui si vuole eseguire la misurazione.\\
	È composto da:
	\begin{itemize}
		\item Un Amperometro
		\item Un Generatore Ideale di Tensione
		\item Una Resistenza Incognita
	\end{itemize}
	\subsubsection{Wattmetro}
	Strumento utilizzato per misurare la potenza di un circuito.\\
	È composto da:
	\begin{itemize}
		\item Un Amperometro
		\item Un Voltmetro
	\end{itemize}
	\subsection{Generatori Pilotati}
	\subsubsection{Generatore di Tensione Pilotato da Una Tensione}
	Legge costitutiva:
	\begin{equation}
		V_{AB} = \alpha \cdot V_k
	\end{equation}
	\begin{itemize}
		\item $\alpha$ è una costante caratteristica del generatore
		\item $V_k$ è la tensione prelevata in un altro punto del circuito
	\end{itemize}
	\subsubsection{Generatore di Tensione Pilotato da Una Corrente}
	Legge costitutiva:
	\begin{equation}
		V_{AB} = R \cdot I_k
	\end{equation}
	\begin{itemize}
		\item $R$ è la resistenza del generatore (Costante) $[\Omega]$
		\item $I_k$ è la corrente prelevata in un altro punto del circuito
	\end{itemize}
	\subsubsection{Generatore di Corrente Pilotato da Una Tensione}
	Legge costitutiva:
	\begin{equation}
		I = G \cdot V_k
	\end{equation}
	\begin{itemize}
		\item $G$ è la conduttanza del generatore (Costante) $[\frac{1}{\Omega}]$
		\item $V_k$ è la tensione prelevata in un altro punto del circuito
	\end{itemize}
	\subsubsection{Generatore di Corrente Pilotato da Una Corrente}
	Legge costitutiva:
	\begin{equation}
		I = \beta \cdot I_k
	\end{equation}
	\begin{itemize}
	\item $\beta$ è una costante caratteristica del generatore
	\item $I_k$ è la corrente prelevata in un altro punto del circuito
	\end{itemize}
	\subsection{Teorema di Millman}
	\begin{theorem}
		In un circuito composto da $n$ rami in parallelo qualsiasi, la tensione di esso è data dal rapporto di tutte le correnti di corto-circuito entranti nel nodo $A$ e tutte le conduttanze incrementali.
		\begin{equation}
			V_{AB} = \frac{\sum_{RAMI} I_{CC\textbf{ (Entranti in A)}}}{\sum_{RAMI} G_{Inc}}
		\end{equation}
		\begin{itemize}
			\item Il numeratore corrisponde alla somma della corrente di ciascun ramo, calcolata come se il ramo stesso fosse staccato dal circuito e collegato in cortocircuito
			\item Il denominatore corrisponde alla somma dell’inverso delle resistenze di ciascun ramo, calcolata come se il ramo stesso fosse staccato dal circuito e collegato in cortocircuito
		\end{itemize}
	\end{theorem}
	Nel caso in cui vi siano $n$ generatori reali di tensione collegati in parallelo ciascuno dei quali è composto da una resistenza $R_i$ e un generatore di tensione $E_i$, allora la formula diventa:
	\begin{equation}
		V_{AB} = \frac{\sum_{i=1}^{n} \frac{E_i}{R_i}}{\sum_{i=1}^{n} \frac{1}{R_i}}
	\end{equation}
	Nel caso il circuito contenga un generatore pilotato che preleva una corrente $I_k$ o una tensione $V_k$ incognita, per calcolare $V_{AB}$ è necessario ricavare tale incognita in funzione di $V_{AB}$.
	\subsection{Trasformazioni Stella-Triangolo}
	Le trasformazioni stella-triangolo permettono di trovare i valori di resistenza dei resistori quando si trovano in una disposizione a stella o a triangolo.\\
	Per prima cosa bisogna trovare le rispettive leggi costitutive quando si trovano a triangolo e a stella:
	\subsubsection{Leggi Costitutive a Stella}
	\[
	\rotatebox[origin=c]{180}{Y} =
	\begin{cases}
		R_{Eq \;\; AB} = R_A + R_B \\
		R_{Eq \;\; BC} = R_B + R_C \\
		R_{Eq \;\; AC} = R_A + R_C
	\end{cases}
	\]
	\subsubsection{Leggi Costitutive a Triangolo}
	\[
	\Delta =
	\begin{cases}
		R_{Eq \;\; AB} = R_C // (R_A + R_B)  = \frac{R_C \cdot (R_B + R_A)}{R_C + R_B + R_A}\\\\
		R_{Eq \;\; BC} = \frac{R_A \cdot (R_B + R_C)}{R_C + R_B + R_A} \\\\
		R_{Eq \;\; AC} = \frac{R_B \cdot (R_C + R_A)}{R_C + R_B + R_A}
	\end{cases}
	\]
	Una volta ottenuti i sistemi di equazioni costitutive possiamo eseguire le trasformazioni semplicemente sostituendo i valori di una nell'altra e viceversa e saltando alcuni passaggi otteniamo:
	\subsubsection{Trasformazione da Stella a Triangolo}
	\begin{equation}
		\rotatebox[origin=c]{180}{Y} \rightarrow \Delta =
		\begin{cases}
			R_A = \frac{R_A \cdot R_B + R_B \cdot R_C + R_C \cdot R_A}{R_A} \\\\
			R_B = \frac{R_A \cdot R_B + R_B \cdot R_C + R_C \cdot R_A}{R_B}\\\\
			R_C = \frac{R_A \cdot R_B + R_B \cdot R_C + R_C \cdot R_A}{R_C}
		\end{cases}
	\end{equation}
	\subsubsection{Trasformazione da Triangolo a Stella}
	\begin{equation}
		\Delta \rightarrow \rotatebox[origin=c]{180}{Y} =
		\begin{cases}
			R_A = \frac{R_B \cdot R_C}{R_A + R_B + R_C} \\\\
			R_B = \frac{R_A \cdot R_C}{R_A + R_B + R_C} \\\\
			R_C = \frac{R_B \cdot R_A}{R_A + R_B + R_C}
		\end{cases}
	\end{equation}
	\subsection{Linearità}
	Un circuito si dice lineare quando i grafici dei suoi segnali d’ingresso e d’uscita sono delle rette, ovvero, dati gli ingressi $x_1$ e $x_2$ e le rispettive uscite $y_1 = f(x_1)$ e $y_2 = f(x_2)$, possiamo definire il circuito lineare se:
	\[f(\alpha x_1 + \beta x_2) = \alpha y_1 + \beta y_2\]
	\subsection{Principio di Sovrapposizione degli Effetti}
	Se un circuito (lineare) ha più generatori (segnali d’ingresso), si possono ricavare le grandezze elettriche risultanti (segnali
	d’uscita) tramite combinazioni lineari dei primi.\\\\
	Esempio:\\
	In un circuito abbiamo come ingressi un generatore di tensione $E_1$ e un generatore di corrente $I_{g1}$, mentre come uscite ha $V_{AB}$, $I_1$ e $I_2$.\\
	Possiamo quindi calcolare le uscite come:
	\[\begin{cases}
		V_{AB} = a_{11} E_1 + a_{12} I_{g1}\\
		V_{AB} = a_{21} E_1 + a_{32} I_{g1}\\
		V_{AB} = a_{31} E_1 + a_{32} I_{g1}\\
	\end{cases}
	\Leftrightarrow
	\begin{bmatrix}
		V_{AB}\\
		I_1\\
		I_2
	\end{bmatrix}
	=
	A_{3 \times 2} \cdot \begin{bmatrix}
		E_1\\
		I_{g1}
	\end{bmatrix}
	\]
	\subsubsection{Metodo della Sovrapposizione degli Effetti}
	Come conseguenza del principio di sovapposizione degli effetti possiamo ottenere il metodo di sovrapposizione degli effetti, che ci permette di calcolare le uscite di circuiti lineari in 2 semplici step:
	\begin{enumerate}
		\item Si accendono singolarmente a turno le vari sorgenti del circuito ignorando le altre e si calcolano le uscite (\textbf{Uscite Parziali})
		\item Si sommano fra di loro le uscite parziali
	\end{enumerate}
	\subsection{Teorema di Thevenin}
	\begin{theorem}
		Dato un bipolo lineare $X$ collegato ad un qualisasi bipolo $Y$\\
		Esiste sempre un generatore reale di tensione ($V$) equivalente ad $X$
	\end{theorem}
	\begin{proof}[Dimostrazione]
		Per ottenere il generatore reale di tensione equivalente si prosegue nel seguente modo:
		\begin{enumerate}
			\item Suddividere il circuito in una rete $X$ ed una rete $Y$, collegate tra loro nei nodi $A$ e $B$.\\
			Quando è richiesto di calcolare una grandezza incognita, solitamente si pone come rete $Y$ la rete contenente tale grandezza.
			\item Calcolare $R_{Eq}$ della rete $X$.\\
			Per fare ciò, si devono spegnere contemporaneamente tutti i generatori indipendenti contenuti in $X$.
			\item Si calcola $E_{Eq'}$ che corrisponde a $V_{AB_{0'}}$, ovvero si lascia aperta la rete $X$ nei poli $A$ e $B$.
			\item Si uniscono i risultati, ottenendo un generatore reale di tensione $E_{Eq}$ e resistenza $R_{Eq}$.
		\end{enumerate}
	\end{proof}
	Esiste anche una versione più estesa del teorema nel quale $X$ contiene $n$ generatori ideali di tensione e $n$ generatori ideali di corrente, ma non verrà approfondità in questi appunti:
	\begin{equation}
		V_{AB} = \sum_{k=1}^{n} H_k \cdot E_k + \sum_{k=1}^{n} L_k \cdot I_{g_k} + K \cdot I
	\end{equation}
	Dove:
	\begin{itemize}
		\item $H_k$ è un coefficiente
		\item $E_k$ sono i generatori di tensione indipendenti
		\item $L_k$ è un coefficiente
		\item $I_k$ sono i generatori di corrente indipendenti
		\item $K \cdot I$ in fondo alla formula è detto Contributo Esterno
	\end{itemize}
	\subsection{Teorema di Norton}
	\begin{theorem}
		Dato un bipolo lineare $X$ collegato ad un qualisasi bipolo $Y$\\
		Esiste sempre un generatore reale di corrente ($I$) equivalente ad $X$
	\end{theorem}
	\begin{proof}[Dimostrazione]
		Per ottenere il generatore reale di tensione equivalente si prosegue nel seguente modo:
		\begin{enumerate}
			\item Suddividere il circuito in una rete $X$ ed una rete $Y$, collegate tra loro nei nodi $A$ e $B$.\\
			Quando è richiesto di calcolare una grandezza incognita, solitamente si pone come rete $Y$ la rete contenente tale grandezza.
			\item Calcolare $R_{Eq}$ della rete $X$.\\
			Per fare ciò, si devono spegnere contemporaneamente tutti i generatori indipendenti contenuti in $X$.
			\item Si calcola $I_{Eq'}$ che corrisponde a $I_{CC'}$, ovvero si mette la rete $X$ in cortocircuito nei poli $A$ e $B$.
			\item Si uniscono i risultati, ottenendo un generatore reale di corrente $I_{Eq}$ e resistenza $R_{Eq}$.
		\end{enumerate}
	\end{proof}
	Un altro metodo per ricavare l’equivalente generatore reale di corrente è seguire il metodo alternativo enunciato
	per il teorema di Thevenin, sostituendo alla fine il generatore reale di tensione con un generatore reale di
	corrente utilizzando le trasformazioni fra generatori reali:
	\begin{equation}
		I_{Eq} = \frac{E_{Eq}}{R_{Eq}}
	\end{equation}
	\subsection{Thevenin e Norton in Presenza di Generatori Pilotati}
	\subsubsection{Metodo 1}
	Calcolo la $R_{Eq}$ come Thevenin diviso Norton:
	\[R_{Eq} = \frac{E_{Eq}}{I_{Eq}}\]
	\subsubsection{Metodo 2 di Norton}
	Applico alla rete $X$ un generatore ideale di tensione $V^*$ e ricavo $I^* = f(V^*)$
	\subsubsection{Metodo 3 di Thevenin}
	Applico alla rete $X$ un generatore ideale di corrente $I'$ e ricavo $V' = f(I')$
	\subsection{Metodi Risolutivi}
	\subsubsection{Metodo delle Correnti di Maglia}
	Basato sul 2° Principio di Kirchhoff.\\
	\begin{enumerate}
		\item Si individuano le varie maglie (Tutte con lo stesso verso)
		\item Si "creano" delle correnti circolari fittizzie per ogni maglia, decidendo liberamente il verso opportuno
		\item Si crea prima la matrice delle resistenze, nella quale ogni elemento corrisponde alla sommatoria delle resistenze delle maglie (negative o positive a seconda del verso della loro corrente)\\
		Esempio con 3 maglie: \\
		$\begin{bmatrix}
			a_{11} & a_{12} & a_{13}\\
			a_{21} & a_{22} & a_{23}\\
			a_{31} & a_{32} & a_{33}\\
		\end{bmatrix}$\\
		Nella seguente matrice $a_{11}$ sarà la sommatoria delle resistenze presenti nella prima maglia, $a_{12}$ sarà la sommatoria delle resistenze in comune fra la prima e la seconda maglia (se esistono) e $a_{13}$ sarà la sommatoria delle resistenze in comune fra la prima e la terza maglia (se esistono).\\
		La matrice è \textbf{simmetrica}.
		\item Si crea la matrice delle tensioni, composta dalla somma delle tensioni dei vari generatori di tensione in ogni maglia, se presenti\\
		Esempio con 3 maglie:\\
		$\begin{bmatrix}
			e_{11} \\
			e_{21} \\
			e_{31}
		\end{bmatrix}$\\
		$e_{11}$ corrisponde alla somma delle delle tensioni dei generatori nella prima maglia e cosi via.
		\item Infine si risolve il sistema lineare.\\
		Esempio con 3 maglie:\\
		$\begin{bmatrix}
			a_{11} & a_{12} & a_{13}\\
			a_{21} & a_{22} & a_{23}\\
			a_{31} & a_{32} & a_{33}\\
		\end{bmatrix}
		\cdot 
		\begin{bmatrix}
			I_{m1} \\
			I_{m2} \\
			I_{m3} 
		\end{bmatrix}
		=
		\begin{bmatrix}
			e_{11} \\
			e_{21} \\
			e_{31}
		\end{bmatrix}$\\
		Dove $I_{m1}$, $I_{m2}$ e $I_{m3}$ sono le correnti di maglia fittizzie
	\end{enumerate}
	Una volta ottenute le correnti di maglia si possono calcolare le varie correnti utilizzando quelle di maglia.\\
	\href{https://www.youtube.com/watch?v=R8OTpGP6NKw}{Link molto utile}\\\\
	Nel caso in cui in una maglia è presente un generatore di corrente si può escludere quella maglia dal calcolo, poichè la sua corrente sarà uguale a quella del generatore.
	\subsubsection{Metodo dei Potenziali di Nodo}
	Basato sul 2° Principio di Kirchhoff.\\
	\begin{enumerate}
		\item Si individuano i vari nodi
		\item Si sceglie un nodo di riferimento (a terra), mentre tutti gli altri sono indipendenti e vengono considerati nei calcoli
		\item Si crea prima la matrice delle conduttanze, nella quale ogni elemento corrisponde alla sommatoria delle conduttanze entranti nei nodi\\
		Esempio con 3 nodi: \\
		$\begin{bmatrix}
			a_{11} & a_{12} & a_{13}\\
			a_{21} & a_{22} & a_{23}\\
			a_{31} & a_{32} & a_{33}\\
		\end{bmatrix}$\\
		Nella seguente matrice $a_{11}$ sarà la sommatoria delle conduttanze entranti nel primo nodo, $a_{12}$ sarà la sommatoria (cambiata di segno) delle conduttanze in comune fra il primo e il secondo nodo (se esistono) e $a_{13}$ sarà la sommatoria (cambiata di segno) delle conduttanze in comune fra il primo e il terzo nodo (se esistono).\\
		La matrice è \textbf{simmetrica}.
		\item Si crea la matrice delle correnti di corto-circuito entranti nei nodi (tenendo conto del verso)\\
		Esempio con 3 nodi:\\
		$\begin{bmatrix}
			i_{11} \\
			i_{21} \\
			i_{31}
		\end{bmatrix}$\\
		$i_{11}$ corrisponde alla somma delle delle correnti ($\frac{e}{R}$) entranti ed uscenti (a seconda del verso) dal nodo.
		\item Infine si risolve il sistema lineare.\\
		Esempio con 3 nodi:\\
		$\begin{bmatrix}
			a_{11} & a_{12} & a_{13}\\
			a_{21} & a_{22} & a_{23}\\
			a_{31} & a_{32} & a_{33}\\
		\end{bmatrix}
		\cdot 
		\begin{bmatrix}
			V_A \\
			V_B \\
			V_C 
		\end{bmatrix}
		=
		\begin{bmatrix}
			i_{11} \\
			i_{21} \\
			i_{31}
		\end{bmatrix}$\\
		Dove $V_A$, $V_B$ e $V_C$ sono i potenziali dei nodi.
	\end{enumerate}
	\href{https://www.youtube.com/watch?v=bGcHdn-mjVg}{Link molto utile}
	\subsubsection{Casi Degeneri}
	\paragraph{Ramo con $G_{Inc} = 0$}
	Se il circuito contiene un ramo con una conduttanza incrementale $= 0$, vorra dire che la corrispondente resistenza tendera ad infinito ($R_{Inc} \rightarrow \inf$) e che di conseguenza una delle equazioni risultanti dell'utilizzo del metodo delle correnti di maglia sarà irrisolvibile.\\
	Quindi in questo caso:
	\begin{itemize}
		\item Non si potrà usare il metodo delle correnti di maglia
		\item Si potrà usare il metodo dei potenziali di nodo
	\end{itemize}
	\paragraph{Ramo con $R_{Inc} = 0$}
	Se il circuito contiene un ramo con una resistenza incrementale $= 0$, vorra dire che la corrispondente conduttanza tendera ad infinito ($G_{Inc} \rightarrow \inf$) e che di conseguenza una delle equazioni risultanti dell'utilizzo del metodo dei potenziali di nodo sarà irrisolvibile.\\
	Quindi in questo caso:
	\begin{itemize}
		\item Non si potrà usare il metodo dei potenziali di nodo
		\item Si potrà usare il metodo delle correnti di maglia
	\end{itemize}
\end{document}